\chapter{Una soglia abitabile: forma, corpi e spazio pubblico}\label{sec:design}

Il Portale porta una domanda di progetto nella vita ordinaria della città:
quanta architettura serve perché una persona si senta autorizzata a cominciare?
Una tastiera disponibile non basta. Occorrono un posto dove fermarsi, un gesto
comprensibile, la libertà di ascoltare senza esporsi e la possibilità che un
altro si avvicini. La forma a goccia del documento del 2016 riunisce struttura,
schermo e diffusione sonora attorno a questo incontro
\cite[pp.~2--4]{terzi_peeragogy_2016}. La sua forza concettuale risiede nella
scala intermedia tra strumento e stanza: abbastanza riconoscibile da diventare
un riferimento, abbastanza raccolta da ospitare un tentativo personale.
Questa lettura orienta il progetto contemporaneo; non è un effetto percettivo
già dimostrato sul pubblico.

\section{La goccia come invito e come limite}

La curva suggerisce una cornice continua attorno al gesto musicale. Può
rendere leggibile, anche da lontano, il rapporto tra chi suona e l'oggetto;
la tastiera offre una presa concreta a una figura altrimenti scultorea.
Il disegno conserva così una tensione feconda: l'oggetto possiede una propria
presenza, ma la sua composizione si completa quando qualcuno lo abita.
La bellezza cercata è quella di una disponibilità riconoscibile: un pezzo
di città che sembra attendere una partecipazione possibile.

La stessa cornice può però produrre un effetto opposto. Una persona può
leggerla come palco, vetrina o apparecchio riservato agli esperti. Occorre
quindi mettere alla prova almeno due orientamenti della postazione: uno
che espone l'esecutore al passaggio e uno obliquo che protegge il primo
tentativo lasciando visibile l'attività. Un prototipo a grandezza reale,
anche privo di elettronica, consentirebbe di osservare esitazioni,
aggiramenti e punti di arresto. Chiedere successivamente che cosa fosse
sembrato permesso aiuterebbe a distinguere curiosità, incomprensione e
timore di disturbare. L'attrattività non coincide con il numero di sguardi:
un oggetto molto fotografato può restare difficile da usare.

\section{Tre geografie della partecipazione}

La configurazione individuale richiede un ingresso semplice, comandi
raggiungibili e la possibilità di terminare senza negoziare con un gruppo.
La seduta deve poter essere spostata da un operatore per liberare
l'avvicinamento; il supporto della tastiera e il volume sottostante
vanno verificati con corpi, posture e ausili differenti. Nessuna quota
ergonomica viene dedotta dall'illustrazione storica.

Con due persone cambia il centro dell'esperienza. Il secondo partecipante
deve vedere le mani, ascoltare e intervenire senza stare necessariamente
alle spalle del primo. Si propone una posizione laterale per il compagno,
una visualizzazione leggibile da entrambi e un accesso alternato agli
strumenti. La condivisione della tastiera può essere un'attività; non deve
essere l'unica modalità di cooperazione. Uno può scegliere il ritmo e
l'altro la frase, oppure ascoltare e restituire un'indicazione. La forma
deve consentire questi scambi senza imporre la postura docente--allievo.

Un piccolo gruppo richiede invece uno spazio di attesa distinto da quello
di esecuzione, turni riconoscibili e un facilitatore nelle sessioni previste.
L'anello degli osservatori non deve chiudere l'uscita né trasformare chi
prova in un concorrente. Non viene fissata una capienza: dipenderà dal
rilievo, dai percorsi e dalle condizioni autorizzate. La prova spaziale
deve includere anche chi attraversa il luogo senza partecipare. Un'attività
civica riuscita lascia praticabile questa scelta.

\section{L'inattività è uno stato del progetto}

La fonte propone la conversione della tastiera in mensola e dello schermo
in galleria di immagini e video \cite[p.~2]{terzi_peeragogy_2016}.
La proposta contemporanea distingue quattro stati percepibili dal pubblico:
disponibile, in uso, in pausa e fuori servizio. Sono la rappresentazione
semplificata degli stati tecnici descritti nel capitolo seguente; la pausa
resta una condizione della sessione in uso. Una galleria silenziosa
può accompagnare la disponibilità, purché un messaggio stabile spieghi
come iniziare. In pausa la persona deve poter riprendere o concludere;
fuori servizio l'oggetto deve dichiararlo chiaramente, senza simulare
un'esperienza ancora funzionante.

La conversione fisica in mensola richiede un dettaglio costruttivo
dedicato: protezione dei tasti, arresto della parte mobile, verifica
degli interstizi e valutazione dei carichi prevedibili. Se tale meccanismo
non è risolto, il prototipo deve limitarsi a una copertura rimovibile
gestita dal personale, senza presentarla come soluzione definitiva.
Con la copertura in posizione l'oggetto è chiuso all'uso musicale:
non deve mostrare lo stato disponibile. L'avvio immediato al primo tasto
vale soltanto dopo la riapertura e il controllo dell'operatore.
L'uso come appoggio rende prevedibili oggetti e bevande: la separazione
dall'elettronica deve entrare nella progettazione fin dall'inizio.
Anche la quiete visiva conta. Una galleria non ha bisogno di lampeggiamenti,
pubblicità o animazione continua per appartenere al giardino.

\section{Accessibilità: una catena, dal percorso alla conclusione}

L'accessibilità comprende il percorso fino alla postazione, l'accostamento,
la postura, l'uso dei comandi, la comprensione dell'attività e il modo di
uscirne. Un'interfaccia leggibile dentro un oggetto irraggiungibile non
risolve la catena. Il sopralluogo deve documentare pavimentazione,
pendenze, soglie, luce utile dei passaggi e spazi occupati da arredi
o attività preesistenti. Le soluzioni andranno provate con persone con
esigenze motorie e sensoriali differenti, riconoscendo limiti e priorità.

Per l'interfaccia digitale, WCAG~2.2 offre criteri verificabili riguardo
a contrasto, uso da tastiera, visibilità del fuoco, alternative alle
informazioni esclusivamente visive e dimensione dei bersagli
\cite{des_wcag22}. Questi criteri costituiscono un riferimento di
progettazione del software; non attestano l'accessibilità fisica del
manufatto. Il gruppo di progetto deve registrare quali criteri verifica
e con quale procedura, evitando dichiarazioni generiche di conformità.

Le attività musicali dovrebbero consentire una prima esplorazione con
una mano e senza pedaliera. Una funzione equivalente azionabile a mano
può rendere disponibili effetti altrimenti legati ai pedali. Ritmo e
turno possono essere indicati anche visivamente; istruzioni e stato
devono restare comprensibili senza affidarsi al solo colore. Queste
alternative non rendono identica ogni esperienza sensoriale: ampliano
le possibilità di partecipazione e richiedono valutazione con gli utenti.
Un comando riconoscibile per fermare suono e attività deve essere sempre
disponibile. L'assistenza umana va resa individuabile senza imporla a chi
può e desidera procedere autonomamente.

\section{L'ascolto condiviso e il diritto al raccoglimento}

La scelta acustica determina quale comunità può formarsi attorno
all'oggetto. Si propongono tre modalità: cuffia individuale, due uscite
per un ascolto condiviso e diffusione locale durante attività concordate.
La prima protegge il tentativo; la seconda rende possibile confrontarsi
sullo stesso materiale; la terza permette un incontro aperto. Il passaggio
alla diffusione deve essere esplicito, con livello iniziale controllato,
senza trasferire improvvisamente al parco un'esecuzione iniziata in privato.

Il progetto acustico richiede misure del rumore di fondo e prove alle
posizioni dell'esecutore, del compagno e dei ricettori esterni pertinenti.
Vanno documentati condizioni delle aperture, attività concomitanti,
posizione dei diffusori e impostazioni. La guida WHO per luoghi ed eventi
musicali richiama controllo dei livelli, monitoraggio e progettazione
acustica \cite{des_who2022}; non fornisce automaticamente il limite
appropriato a questo giardino. Obiettivi e condizioni applicabili andranno
definiti nel progetto acustico specifico. Una percentuale del cursore
volume non equivale a una misura dell'esposizione all'orecchio.

Per le cuffie servono un insieme controllato di modelli, una procedura
di pulizia e una configurazione verificata dell'uscita. La possibilità
di collegare cuffie personali va valutata rispetto alla variabilità
della loro risposta. Non si assume che la sola cuffia risolva la privacy:
schermo, voce e posizione del corpo possono rendere osservabile la sessione.
Il contenuto personale non deve apparire nella galleria successiva.
La webcam prevista dalla fonte deve poter restare fisicamente oscurata;
una sessione remota richiede un'inquadratura che non inglobi i passanti.

\section{Luce, clima e cura materiale}

Nel Padiglione candidato il rapporto tra vetro e giardino merita una
verifica sul posto: riflessi, contrasto dello schermo, ombre sulle mani
e visibilità dall'esterno cambiano con ora e stagione. Un render con
illuminazione favorevole non dimostra la leggibilità reale. Si propone
di confrontare orientamenti dello schermo e finiture campione, verificando
le schermate operative in più condizioni di luce. Una soluzione che
richieda oscuramenti permanenti deve essere discussa anche per il suo
effetto sull'architettura e sulle altre attività.

Temperatura, umidità, condensa e ventilazione richiedono una campagna
di misura e il confronto con le specifiche dei componenti scelti.
L'installazione interna e quella temporanea all'aperto sono configurazioni
diverse: un involucro adatto alla prima non viene dichiarato adatto alla
seconda. Calore prodotto dal computer e rumore delle ventole possono
compromettere il raccoglimento. La manutenzione deve raggiungere filtri,
connettori e parti sostituibili senza smontare l'intera scultura.

Stabilità, spigoli, parti mobili, trazione dei cavi e accesso ai componenti
elettrici costituiscono verifiche del progetto esecutivo, affidate a
competenze appropriate. La gestione quotidiana deve includere controllo
visivo, pulizia compatibile con i materiali, segnalazione dei guasti
e messa fuori uso dei componenti danneggiati. Il rapporto con un bene
pubblico si costruisce anche qui: una riparazione tempestiva comunica
che l'invito rivolto ai cittadini è ancora mantenuto.

\section{Reversibilità e modello verificabile}

Si propone una prima configurazione smontabile, con protezioni nei
punti di appoggio e percorsi dei cavi rimovibili. L'assenza di fissaggi
non garantisce da sola stabilità né autorizzazione: massa, baricentro,
carichi sul pavimento, movimenti prevedibili e modalità di trasporto
vanno valutati insieme. Ogni intervento sull'edificio resta subordinato
alle verifiche e alle autorizzazioni pertinenti. Il criterio di
reversibilità comprende ripristino del luogo, tempo di rimozione e
possibilità concreta di continuare le attività già presenti.

Il modello Blender deve servire a confrontare ipotesi e a rendere
leggibile la provenienza dei dati. Per la versione di ricerca si propone
di distinguere oggetti \texttt{\_PROXY}, dimensioni stimate
\texttt{\_EST} e geometrie rilevate corredate da data e metodo.
Un piano di riferimento comune, unità esplicite e parametri separati
per involucro, tastiera, seduta, schermo e ingombri di utilizzo rendono
possibili confronti riproducibili. Le figure umane sono involucri di
prova, non sostituti di una valutazione ergonomica. Le visualizzazioni
devono riportare in didascalia lo stato del rilievo; il termine
\emph{gemello digitale} sarebbe prematuro senza una corrispondenza
verificata e un processo di aggiornamento.

\begin{table}[htbp]
\centering\small
\caption{Dati necessari e prove proposte. Le soglie di accettazione fisica
e acustica saranno definite prima della sperimentazione, sul sito e sulla
configurazione effettivi.}
\label{tab:design-prove}
\begin{tabular}{P{0.23\linewidth}P{0.32\linewidth}P{0.35\linewidth}}
\hline
Parametro da acquisire & Metodo e provenienza & Criterio di prova \\
\hline
Percorsi e aperture & Rilievo quotato; fotografie datate & Avvicinamento, uscita e transito provati nelle tre configurazioni \\
Tastiera e comandi & Quote del prototipo; specifiche dei componenti & Attività iniziale completabile nelle posture considerate; ostacoli registrati \\
Schermo e riflessi & Prove in sito con orario e luce documentati & Istruzioni leggibili dai posti previsti, senza dover cambiare postura \\
Emissione sonora & Misure con strumenti e configurazione documentati & Rispetto degli obiettivi definiti; passaggio tra modalità senza picchi inattesi \\
Clima e dissipazione & Serie temporali; dati dei produttori & Condizioni entro gli intervalli ammessi; arresto gestito fuori intervallo \\
Stabilità e appoggi & Calcolo e prove del progetto esecutivo & Esito positivo delle verifiche previste prima dell'uso pubblico \\
Rimozione e manutenzione & Prova completa, tempi e utensili registrati & Ripristino documentato; parti di consumo raggiungibili e sostituibili \\
\hline
\end{tabular}
\end{table}

La sequenza proposta procede dal modello volumetrico al campione
ergonomico a grandezza reale, quindi al prototipo funzionante e alla
prova presidiata in sito. Ogni passaggio deve poter cambiare il disegno.
Se per conservare la silhouette si sacrifica la possibilità di avvicinarsi,
la forma ha perso il proprio compito. L'ambizione del Portale è tenere
insieme riconoscibilità e ospitalità: fare della cura del dettaglio la
condizione concreta di un incontro che prima non aveva luogo.
